\section{The birational threshold reduction}
\label{sec:reduction}

The central assertion is the following uniform estimate.

\begin{proposition}\label{prop:birational-threshold}
For every positive integer $d$ and every $\epsilon>0$ there is
$c(d,\epsilon)>0$ such that, over $\C$, the following holds. Let
$f\colon X\to Z$ be a birational contraction of $d$-dimensional normal
varieties, let $X$ be of Fano type over $Z$, and assume that $K_X$ is
$\Q$-Cartier, $-K_X$ is nef over $Z$, and $X$ is $\epsilon$-lc. For each
closed point $z\in Z$,
\begin{equation}\label{eq:A1}
 \lct\bigl(X;\m_z\OO_X\bigr)\ge c(d,\epsilon).
\end{equation}
\end{proposition}

The subsequent sections prove Proposition~\ref{prop:birational-threshold}.
We first explain precisely why it suffices.

\begin{lemma}\label{lem:threshold-to-complement}
Proposition~\ref{prop:birational-threshold} in all dimensions implies
Theorem~\ref{thm:main} over $\C$.
\end{lemma}

\begin{proof}
Shrink $Z$ to an affine neighbourhood of $z$. Choose generators of $\m_z$
and take a log resolution of their pullback ideal and of $X$. The pullback
system has a fixed divisor and a base-point-free moving system. By
\eqref{eq:A1}, multiplying the fixed divisor by $c(d,\epsilon)$ preserves
log canonicity. For a general linear combination $h$ of the generators,
Bertini's theorem makes its moving divisor transverse to the resolution
strata. Consequently, with
$t=\min\{c(d,\epsilon),1/2\}$, the pair
$(X,t f^*\Div(h))$ is lc near $f^{-1}(z)$; the same conclusion follows
componentwise on the resolution. The divisor $\Div(h)$ passes through $z$.
This is the birational boundary-zero instance of the local hypersurface
assertion in \cite[Conjecture~1.3]{Chen}.

Apply the implication (4)$\Rightarrow$(1) in
\cite[Theorem~1.6]{Chen}, for every base-dimension bound needed and with
boundary coefficient set $\{0\}$. Its conclusion supplies a bounded klt
complement on an $\epsilon$-lc contraction of Fano type whose
anti-log-canonical divisor is relatively nef. A Fano contraction as in
Theorem~\ref{thm:main} satisfies these hypotheses, using the zero boundary.
Over our affine base its total space is quasi-projective, as required there.
Relative linear triviality becomes actual linear triviality after shrinking
about $z$ to trivialize the line bundle on the base.

The birational assertion in Chen's formulation allows a canonical Weil
divisor that is $\R$-Cartier. Such a divisor is already $\Q$-Cartier: on a
finite open cover, expressing its Cartier condition in terms of finitely many
integral divisors gives a finite system of rational linear equations. A real
solution of this consistent system can be replaced by a rational solution,
and a common denominator works on the finite cover. Finally a real positive
lower discrepancy bound may be decreased to a positive rational one. Thus
these harmless differences in formulation do not restrict the implication.
\end{proof}

\begin{lemma}\label{lem:field-descent}
If Theorem~\ref{thm:main} holds over $\C$ with a fixed index depending only
on $d$ and $\epsilon$, it holds over every algebraically closed field of
characteristic zero with a fixed multiple of that index.
\end{lemma}

\begin{proof}
Work over an affine neighbourhood of $z$. Descend the morphism, the point,
and the relevant canonical-divisor data to a countable algebraically closed
subfield $k_0$ of the given field $k$. All defining equations require only
finitely many coefficients, so such a $k_0$ exists and embeds in $\C$.
Normality, the contraction property, and relative ampleness descend and
ascend under these field extensions after enlarging the defining field when
necessary. A log resolution defined over $k_0$ checks the discrepancy bound:
its finitely many boundary coefficients, together with the smooth transverse
strata, compute the same lc condition after either extension. We may
therefore apply the complex assertion to this descended model.

Let $b>1$ be a fixed multiple of the complex complement index. Write $K_X$
using a rational top-degree form $\omega$ defined over $k_0$. The coherent
module of relative sections of the reflexive sheaf $\OO_X(-bK_X)$ has
finitely many generators over the affine base. Represent these generators
by rational functions $\phi_1,\ldots,\phi_s$ over $k_0$. Sections commute
with flat constant-field extension; equivalently one may compute the
reflexive powers on the smooth locus and use flat base change there.
We do not require $bK_X$ to be Cartier.

Take a log resolution $\pi\colon Y\to X$ over $k_0$ on which the fractional ideal
generated by the $\phi_i$ is principal, and its divisor together with
$\Div_Y(\omega)$ and the exceptional locus has simple normal crossings.
For every component $E$ of these divisors put
\[
 a_E=\ord_E(\omega),\qquad
 q_E=\min_i E(\phi_i).
\]
Over $\C$, the given klt complement is represented, after localizing on a
neighbourhood of $z$, by a section $\phi$ with
$\Div_X(\phi)=b(K_X+B)$. Its order at every $E$ whose proper image in the
base contains $z$ is at least $q_E$: the coefficients expressing it in the
chosen generators are regular on that neighbourhood. Klt of this complement
therefore gives the strict inequalities
\begin{equation}\label{eq:descent-fixed-part}
 1+a_E>\frac{q_E}{b}
\end{equation}
for all such components. These are finite inequalities for the descended
resolution, so they remain true over $k$.

Remove from the base the proper images of the finitely many offending
components that do not contain $z$. The fractional ideal is now its fixed
divisor times a base-point-free system. A general $k$-linear combination
$\phi=\sum_i\lambda_i\phi_i$ has a moving divisor which is smooth and
transverse to the resolution strata, by Bertini in characteristic zero.
Its coefficient in
\[
 \frac{1}{b}\Div_Y(\phi)-\Div_Y(\omega)
\]
is $1/b<1$; the fixed coefficients are strictly less than one by
\eqref{eq:descent-fixed-part}. This simple-normal-crossing subpair is klt.
Pushing down gives an effective boundary
\[
 B=\frac{\Div_X(\phi)-bK_X}{b},
 \qquad b(K_X+B)=\Div_X(\phi),
\]
which is klt on a neighbourhood of $f^{-1}(z)$. Properness justifies the
shrinking used above. Closing up this effective divisor in the original $X$
gives the divisor required by the theorem; its restriction to that
neighbourhood is unchanged. Thus the fixed index $b$ works over $k$.
\end{proof}

\subsection{The contradiction sequence}
\label{subsec:contradiction-sequence}

Suppose Proposition~\ref{prop:birational-threshold} fails for fixed
$d>0$ and $\epsilon>0$. Choose a sequence of its data with threshold tending
to zero. We will usually suppress the sequence index. Before constructing
the cone, we make two reductions that preserve this failure.

\begin{lemma}\label{lem:ample-model}
The sequence may be chosen with $-K_X$ ample over $Z$. This replacement
preserves all discrepancies and the threshold in \eqref{eq:A1}.
\end{lemma}

\begin{proof}
Choose a klt log Fano boundary $\Delta$ over $Z$ and a positive integer $r$
with $-rK_X$ Cartier. The divisor $D=-rK_X$ is relatively nef, and for a
positive integer $a$,
\[
 aD-(K_X+\Delta)=-arK_X-(K_X+\Delta)
\]
is relatively ample. The relative klt base-point-free theorem
\cite[Theorem~1.3]{FujinoBPF} makes $-K_X$ semiample over $Z$. Let
$g\colon X\to X^{\mathrm a}$ be its relative ample model. Since $X\to Z$
is birational, so is $g$. Choose canonical divisors from the same rational top form. For a Cartier divisor $H$ on $X^{\mathrm a}$ ample over $Z$ and a sufficiently divisible $m$, write
$-mK_X=g^*H+\Div_X(\phi)$. Birationality identifies the function fields, so
pushing forward gives $-mK_{X^{\mathrm a}}=H+\Div_{X^{\mathrm a}}(\phi)$.
Pulling back this identity gives
\[
 K_X=g^*K_{X^{\mathrm a}}.
\]
The morphism is crepant; the new variety is $\epsilon$-lc and has ample
anti-canonical divisor over $Z$. In particular it is of Fano type there.
Every divisorial valuation has the same discrepancy on these two models,
and its value on the ideal from $Z$ is the same. The valuative formula for
the threshold proves the assertion.
\end{proof}

We henceforth assume that $Z$ is affine and $-K_X$ is ample over $Z$.
By the valuative description of the failing threshold, there are prime
divisors $P$ over $X$ such that
\begin{equation}\label{eq:bad-prime-ratio}
 \frac{P(\m_z)}{A_X(P)}\longrightarrow\infty.
\end{equation}
Here and below $P(\m_z)$ denotes the order of $\m_z\OO_X$.

Birkar's relative bounded lc complement theorem
\cite[Theorem~1.8]{BirkarComplements} applies to the pair $(X,0)$: the
coefficient set is fixed, the total space is of Fano type over the base,
and $-K_X$ is relatively nef. Our base has positive dimension $d$, so the
stated relative theorem applies. After shrinking about $z$, it supplies an
lc complement of an index bounded in terms of $d$. Taking a fixed multiple
of the finitely bounded indices and then a further multiple if necessary,
we obtain one integer $n>1$, independent of the sequence, for which
\begin{equation}\label{eq:auxiliary-lc-system}
 \frac{1}{n}|-nK_X|\quad\text{is an lc linear system}.
\end{equation}
This means that for every divisorial valuation $E$ with centre in the
neighbourhood under consideration,
$\frac1n E(|-nK_X|)\le A_X(E)$, where the system order is the minimum
section order. One lc section already proves these inequalities for the
whole system. The integer $n$ will remain fixed in the cone construction;
it is not asserted to be a klt complement index.

\subsection{The finite-coefficient consequence}

The crepant model also gives the consequence stated in the introduction,
once Theorem~\ref{thm:main} has been proved.

\begin{proof}[Proof of Corollary~\ref{cor:finite-coefficients}]
We check item (2) of \cite[Theorem~1.6]{Chen} in every birational dimension.
For its boundary-zero, relatively nef data, use the crepant ample model
$g\colon X\to Y$ of Lemma~\ref{lem:ample-model}. Apply
Theorem~\ref{thm:main} to $Y\to Z$, decreasing $\epsilon$ to a fixed
positive rational number when needed. After shrinking about $z$, its
complement boundary $B_Y$ is $\Q$-Cartier, since both $K_Y$ and
$K_Y+B_Y$ are $\Q$-Cartier there. On this neighbourhood,
$B_X=g^*B_Y$ is effective and
\[
 K_X+B_X=g^*(K_Y+B_Y).
\]
The pair upstairs is klt, and the same complement index makes its adjoint
divisor Cartier and trivial. Take the effective closure of $B_X$ outside
this neighbourhood if necessary. In particular the multiple used to construct
the ample model does not change the complement index. This proves Chen's
item (2) for every base-dimension bound. The implication (2)$\Rightarrow$(1)
of his theorem gives the assertion, with the stated dependence on $I$.
\end{proof}

In the rest of the paper we work with the contradiction sequence above,
with its fixed auxiliary lc degree $n$, and prove
Proposition~\ref{prop:birational-threshold}.
